\section{Conviction and Rhetoric}

\begin{quotex}
That henceforth we be no more children tossed to and fro, and carried about with every wind of doctrine by the wickedness of men, by cunning craftiness, by which they lie in wait to deceive.
\flright{\textsc{Ephesians 4:14}}
\end{quotex}

Rather than cluttering up the section on \textbf{Carlo Michelstaedter}\footnote{\url{See Section \ref{sec:Michelstaedter} in this book.}} from \textit{Essays on Magical Idealism} with a commentary, let this serve as an introduction. Yet nothing, or almost nothing, is posted on Gornahoor without a reason, without being part of a larger undertaking, without being related to what has been written before.

In his philosophical writings, we see that \textbf{Julius Evola} had developed a complex world view before he was 30. Since the majority of readers here are men under 30, take this as a challenge to do the same. Create a redoubt in the mind. It will protect you from the winds of doctrine. Even if, for a time, it must be placed aside, it will return to you when you need it.

The bulk of men, even those who consider themselves educated and intellectual, are content to wander through the orchard of ideas, picking off the lowest and prettiest fruits, accumulating them in a basket. The wiser, however, will go for the higher and harder-to-get fruit, while the wisest of all, will go to the trunk and look for the source in the roots.

Thus, Julius Evola announces his plan and his indebtedness to certain predecessors. He sees himself as their fulfillment, drawing out unexpected conclusions, and amplifying certain points. Sometimes, a misreading is necessary to move the argument forward. The first figure Evola points to is \textbf{Carlo Michelstaedter}, specifically to his book Persuasion and Rhetoric. By now, his story should be known. He finished his dissertation and then immediately shot himself dead at the age of 23. Presumably, there was a psychological problem and his philosophy does not entail such a radical end to life.

In a few dialogs, Plato distinguished ``rhetoric and persuasion'' from ``wisdom and reason''. The city, as it is, is governed by rhetoric and persuasion. In other words, rhetorical techniques, primarily emotionally based, are used to persuade men to a belief on a course of action. Plato claimed instead that the Republic should be guided by wisdom and reason. Now this should be common knowledge among all educated men today, but it certainly is not. Even worse, in our time, there are no wise men, and reason cannot persuade anyone, since few are capable of reason.

So even the most intelligent use their gifts as rhetoric. They are content to display the fruit they gathered, perhaps creating a sweet tasting concoction out of it. Even in our world of Tradition, rhetoric dominates. That is why we see long quotes from Evola or Nietzsche posted everywhere. Or youtube videos consisting of nothing but quotes, gathered in a heap, with some sort of martial sounding music in the background. In other words, a rhetorical appeal to emotion.

Instead, we must dig into the roots of thought, to read carefully and completely. A good writer will leave nothing in by accident, as he expects to be understood. Men should avail themselves of the opportunities offered them, since they are not offered very often. But the state of nature is sleep, and sleeping men can gather and eat fruit.

\paragraph{Absolute Self}


Michelstaedter, in his misreading, contrasts rhetoric and persuasion. Rhetoric is inauthentic, it means to be immersed in the world of things, in becoming, in giving up oneself to others. Persuasion, on the other hand, is to be in full possession of oneself, to not depend on others, to not be driven by deficiency or lack.

But rhetoric is also speech. So if persuasion is its opposite, then it is mute and dumb. It can only be and cannot express itself. It is absolute and one with all things. This will come out in Evola's commentary. In Italian, it seems to me, persuasion is a stronger word than in English. It can also mean conviction. So I would call the persuaded man, the ``convinced man'', but to do so would introduce an inconsistency.

Of course, the Absolute I or Self is Evola's cardinal point, so Michelstaedter serves as a support for that idea. Also for Jung, whom Evola ignores, the Self must be created since it is not part of nature. Jung also recognizes that Christ serves as the idea of the Self, although that is largely forgotten. The Logos comes in our second birth, not given to all, since it is not part of our natural birth, which is given to all.

\paragraph{Potentiality}


Another common rhetorical device is the claim that we all have ``infinite potential'' or something similar. This makes people exult in their own unfulfilled promise; to deny that idea is considered heresy. Yet what distinguishes men is not their potentiality but their actuality. The measure of power is what they can bring into manifestation. Of course, we must emphasize that the proper measure is the difficulty of the thing, since most things are rather easy to manifest.

In Michelstaedter's philosophy, mere potentiality is simply non-being, a lack, a deficiency. As such it belongs to rhetoric. Sure, we can talk about it, since that makes us feel so good, but it is still far from persuasion.

\paragraph{The Bodhisattva and the Jivan Mukti}


This is still the difference between a Guenon and an Evola. The former aims to be a Jivan Mukti, to transcend the world, into the Absolute. Apparently, a jivan mukti is indistinguishable from a dead man\footnote{\url{http://www.telegraph.co.uk/news/worldnews/asia/india/10860998/Indian-court-asked-to-rule-on-whether-Hindu-guru-dead-or-meditating.html}}.

The latter aims for the Absolute Self, to be self-sufficient. Like a bodhisattva, the enlightened one does not disappear, but returns to show others the way.


\flrightit{Posted on 2014-05-30 by Cologero}

\begin{center}* * *\end{center}

\begin{footnotesize}
\begin{sffamily}
\texttt{Synodius on 2014-06-05 at 15:58 said:}

I don't quite understand that distinction between Guenon and Evola.

In his essay \textit{Ascending and Descending} Realization Guenon writes:

``... the difference in question is finally that which exists between the Pratyeka-Buddha and the Bodhisattva. In this regard, it is particularly important to note that the way which has the first of these two states as its end is designated as a `small way' or, if one wishes, a `lesser way\textquotesingle(hinayana), thus implying that it is not exempt from a certain restrictive character, whereas the one leading to the second state is considered to be truly the `greater way', and therefore the one that is complete and perfect in all respects. This allows us to answer the objection which could arise from the fact that in a general way the state of Buddha is regarded as superior to that of Bodhisattva. In the case of the Pratyeka-Buddha, that superiority can only be apparent, and it is due above all to the character of `impassiveness' which, apparently, the Bodhisattva does not have; we say `apparently' because it is necessary to differentiate between the `reality' of the being and the role it has to play with respect to the manifested world, or in other words, between what it is in itself and what it appears to be for ordinary beings.''

Is there any essential difference between Guenon and Evola on this question?

\hfill

\texttt{Cologero on 2014-06-06 at 07:03 said:}

Although Evola has indicated where he thinks Guenon had gone wrong, I am relying more on impressions.

Guenon distinguishes between what he labels ``salvation'' and ``liberation''. The former is the fulfillment of the human condition while the latter is the fulfillment of all possible conditions.

These states are perhaps clearer in some other authors in our pantheon. E.g., Tomberg represents these two states in the \textit{Emperor}\footnote{\url{http://www.meditationsonthetarot.com/alchemical-transformation}} and the Pope. The Emperor is the ideal human being and has authority (we cannot here refer to other traditions such as the Sufi man of light, etc).

In our \textit{Monday study group}\footnote{\url{http://www.gornahoor.net/?p=7205}}, we again see Mouravieff make the distinction between two higher states of consciousness: (1) Consciousness of the Real I and (2) Absolute consciousness.

I don't see Evola making that distinction and it seems to me that he even confuses them. Yet, he is excellent in describing the consciousness of the Real I. That is why we have made use of his writings so often, since that is probably the most that can be described over the Internet.

But this leads to two limitations, although Evola does not see them as such. On the contrary, he has criticized Guenon for not recognizing them.

The first is his insistence on ``action'' as a path on a par with ``contemplation''. That is true as far as reaching the first higher state, but not for the second.

The second is his insistence that ultimate authority resides in the Emperor, not in a separate spiritual or religious authority. Again, that is true on the human plane, but not for the divine plane. We have been dealing with the Emperor on the \textit{Medtarot Discussion List}\footnote{\url{http://www.meditationsonthetarot.com/medtarot-discussion-list}} and will be discussing the role of the spiritual authority over the next several weeks.

The second stage of higher consciousness requires a spiritual/religious orientation, which Evola eschews, and even rejects.

\hfill

\texttt{Synodius on 2014-06-06 at 10:47 said:}

Thank you for your clarification, I'll do some cross-reading with this in mind.

\end{sffamily}
\end{footnotesize}

